\documentclass[11pt]{scrartcl}

\usepackage[stixtwo]{fontsetup}
\usepackage{amsmath}

\begin{document}
Az $A = \sqrt[3]{6A^2 - 5A}$ egyenlet megoldásai $A_1 = 0$, $A_2 = 1$ és
$A_3 = 5$.

Teljes indukcióval belátjuk, hogy $1 \le a_n \le 5$ minden $n$-re.

$n = 1$ esetén $1 \le a_1 = 2 \le 5$ teljesül.

Ha $1 \le a_n \le 5$, akkor $1 \le 6a_n - 5 \le 25$, ezért
$1 \le a_n (6a_n - 5) = 6a_n^2 - 5a_n \le 125$, tehát
$1 \le \sqrt[3]{6a_n^2 - 5a_n} = a_{n + 1} \le 5$.

Megmutatjuk, hogy $a_n \le a_{n + 1}$ minden $n$-re.

Köbre emelés után $a_n^3 \le 6a_n^2 - 5a_n$, azaz
$a_n (a_n^2 - 6a_n + 5) = a_n^3 - 6a_n^2 + 5a_n \le 0$, ami teljesül, mivel
$a_n > 0$ és $a_n^2 - 6a_n + 5 \le 0$ az előbbi korlátok között.

Megmutattuk, hogy $a_n$ monoton és korlátos, így a konvergencia elégséges
feltétele miatt konvergens.

Legyen $A = \lim a_n \in \mathbb{R}$! Felhasználva a részsorozat
határértékére és a határérték és a műveletek kapcsolatára vonatkozó
tételeket, kapjuk az $A = \sqrt[3]{6A^2 - 5A}$ egyenletet.

Mivel $a_n$ monoton növő, ezért $A \ge a_1 = 2$, így a fenti gyökök közül
csak az $A_3 = 5$ lehet a határérték.
\end{document}
